\documentclass[11pt,a4paper]{amsart}

\usepackage[margin=1in]{geometry}
\usepackage{amsmath,amssymb,amsthm,mathtools}
\usepackage{microtype}
\usepackage[hidelinks]{hyperref}

\newtheorem{theorem}{Theorem}[section]
\newtheorem{corollary}[theorem]{Corollary}
\newtheorem{lemma}[theorem]{Lemma}
\newtheorem{proposition}[theorem]{Proposition}
\theoremstyle{remark}
\newtheorem{remark}[theorem]{Remark}

\newcommand{\aff}{\operatorname{aff}}
\newcommand{\ri}{\operatorname{ri}}
\newcommand{\ex}{\operatorname{ex}}
\newcommand{\ran}{\operatorname{ran}}
\newcommand{\supp}{\operatorname{supp}}
\newcommand{\Borel}{\mathfrak B}
\newcommand{\calH}{\mathcal H}
\newcommand{\calK}{\mathcal K}
\newcommand{\inner}[2]{\left\langle #1,#2\right\rangle}

\title{Planar Strictly Convex Hyperrigidity}
\author{DannyExperiments}
\thanks{The ordinary proof was generated with OpenAI GPT-5.6 Pro.
Manuscript preparation, audit coordination, and provenance tooling used
OpenAI Codex.  DannyExperiments curated the manuscript and associated
research artifacts.}
\date{July 24, 2026}

\begin{document}

\begin{abstract}
Let \(K\) be a compact convex set of affine dimension at most two and let
\(g\colon K\to\mathbb R\) be continuous and strictly convex.
We prove that a normalized positive-operator-valued measure and a
projection-valued measure on \(K\) are equal whenever they agree on every
affine function and on \(g\).
Equivalently, the function system generated by the affine functions and
\(g\) is hyperrigid in \(C(K)\).
Consequently, equality under compression for one such \(g\) forces the
compression subspace to reduce the entire commuting self-adjoint tuple.
No smoothness of \(g\), finite-rank assumption, separability assumption on
the acting Hilbert space, or atomicity of the spectral measure is required.
\end{abstract}

\maketitle

\section{Introduction}

Let \(K\) be a compact convex subset of a finite-dimensional real affine
space.  Write \(A(K)\) for the real affine functions on \(K\), regarded as a
self-adjoint function system in \(C(K)\).  Given a continuous strictly convex
function \(g\colon K\to\mathbb R\), consider
\[
   S_g=\operatorname{span}_{\mathbb C}\bigl(A(K)\cup\{g\}\bigr).
\]
Scalar strict convexity makes every point of \(K\) a Choquet boundary point
for \(S_g\).  Operator hyperrigidity is stronger: it requires uniqueness not
only for scalar representing measures, but for positive-operator-valued
measures representing arbitrary, possibly diffuse, spectral
representations.

Our main result establishes this operator uniqueness in affine dimension at
most two.

\begin{theorem}\label{thm:main}
Let \(K\) be compact and convex with \(\dim\aff K\leq 2\).  Let \(E\) be a
normalized positive-operator-valued measure on \(K\), and let \(F\) be a
projection-valued measure on \(K\), acting on the same Hilbert space.  Let
\(g\colon K\to\mathbb R\) be continuous and strictly convex.  Suppose
\[
 \int_K\ell\,dE=\int_K\ell\,dF\quad(\ell\in A(K)),
 \qquad
 \int_Kg\,dE=\int_Kg\,dF.
\]
Then \(E=F\).
\end{theorem}

Here and below a positive-operator-valued measure (POVM) is countably
additive in the weak operator topology.  Its scalarizations on a compact
metric space are finite regular Borel measures.  A projection-valued measure
(PVM) is a POVM whose values are orthogonal projections.

The corresponding compression statement answers the motivating
two-variable question affirmatively.

\begin{corollary}\label{cor:square}
Let \(X,Y\in B(\calK)\) be commuting positive contractions, let \(P\) be an
orthogonal projection, and put
\[
 H=\ran P,\qquad A=PXP|_H,\qquad B=PYP|_H.
\]
Assume that \(A\) and \(B\) commute.  If
\(f\colon[0,1]^2\to\mathbb R\) is continuous and strictly convex and
\[
   P f(X,Y)P|_H=f(A,B),
\]
then \(PX=XP\) and \(PY=YP\).
\end{corollary}

The proof has two distinct components.  The first is a local interior
argument.  Supporting-plane defects of \(g\) give nonnegative Bregman
functions.  Douglas factorization would turn any off-diagonal spectral
coupling into a nonzero scalar measure of finite reciprocal-Bregman energy.
Minty's parametrization of the subgradient graph and a critical
inverse-square estimate show that such a measure cannot exist in dimensions
one or two.

The second component treats the boundary.  Flat boundary faces are exposed
affinely and reduced to the one-dimensional theorem.  The remaining extreme
boundary may be curved and may carry diffuse spectral measure.  We replace
each point of \(K\) measurably by an extreme-point representing measure and
use Scherer's planar extreme-boundary rigidity theorem
\cite{SchererPlanar}.  Strict Jensen equality then shows that this
boundaryization moved no mass.  Positivity finally glues the interior, flat
faces, and extreme boundary and annihilates all cross terms.

The dimension-one part of Theorem~\ref{thm:main} is already contained in
Brown's one-variable strict-convexity theorem \cite{BrownConvergence}.
Known genuinely planar special cases include positive-definite quadratic
functions, by Shankar's normal-generator theorem \cite{Shankar}, and certain
explicitly specified monomial systems in a single normal generator covered
by Pietrzycki--Stochel \cite{PietrzyckiStochel}.  Davidson--Kennedy
\cite{DavidsonKennedy} give the general relationship among Choquet order,
unique extension, and hyperrigidity for function systems.

There is a useful convex-geometric reformulation.  Set
\[
 \Gamma_g=\{(x,g(x)):x\in K\},\qquad
 L_g=\operatorname{conv}\Gamma_g.
\]
Restriction to \(\Gamma_g\) identifies \(A(L_g)\) completely order
isomorphically with \(S_g\), and strict Jensen equality gives
\(\ex L_g=\Gamma_g\).  Consequently,
\[
 S_g\text{ is hyperrigid in }C(K)
 \quad\Longleftrightarrow\quad
 A(L_g)\text{ is hyperrigid in }C(\ex L_g).
\]
If \(g\) is nonaffine, then
\(\dim\aff L_g=\dim\aff K+1\).  Thus Scherer's planar theorem already handles
the graph lift when \(\dim\aff K=1\), whereas an original planar \(K\)
produces a three-dimensional graph hull.  The scope not covered by those
results is precisely arbitrary continuous strictly convex \(g\) on such
two-dimensional \(K\), including diffuse representations.  Two separately
run solution-aware literature searches completed on July~24, 2026 located no
prior theorem with that full scope; this negative search is not a claim of
historical priority.

Scherer's later spectrahedra theorem \cite{SchererSpectrahedra} proves
hyperrigidity of \(A(L)\) in \(C(\ex L)\) for compact spectrahedra \(L\)
with closed extreme boundary.  Since \(\ex L_g=\Gamma_g\) is already compact,
it applies here precisely when the graph hull \(L_g\) is a spectrahedron.
An arbitrary graph hull of a continuous strictly convex function need not be
a spectrahedron.  Thus that result is closely related but does not subsume
Theorem~\ref{thm:main}.

\section{Interior localization}

We identify \(\aff K\) with \(\mathbb R^d\), where \(d\leq2\), and take all
subgradients relative to the single common domain \(K\).  For
\(z\in\ri K\), set
\[
 \partial_Kg(z)=
 \left\{p\in\mathbb R^d:
 g(t)\geq g(z)+\inner{p}{t-z}\ \text{for every }t\in K\right\}.
\]

\begin{lemma}[Compact-fiber selection]\label{lem:borel-selection}
Let \(X,Y\) be compact metric spaces and let
\(\mathcal R\subset X\times Y\) be closed with every fiber
\(\mathcal R_x=\{y:(x,y)\in\mathcal R\}\) nonempty.  There is a Borel map
\(s\colon X\to Y\) such that \(s(x)\in\mathcal R_x\) for every \(x\).
\end{lemma}

\begin{proof}
Choose a dense sequence \((q_j)\) in \(Y\) and put
\(d_j(x)=\operatorname{dist}(q_j,\mathcal R_x)\).  Each \(d_j\) is lower
semicontinuous.  Indeed, if \(x_m\to x\), pass to a subsequence realizing
the lower limit and choose minimizers \(y_m\in\mathcal R_{x_m}\).  A further
subsequence has limit \(y\in\mathcal R_x\), because \(\mathcal R\) is
closed, and hence
\[
 d_j(x)\leq d(q_j,y)=\liminf_m d_j(x_m).
\]

Set \(\varepsilon_n=2^{-n}\).  Let \(s_1(x)\) be the first \(q_j\) with
\(d_j(x)<\varepsilon_1\).  Once \(s_n\) is defined, let \(s_{n+1}(x)\) be
the first \(q_j\) satisfying
\[
 d_j(x)<\varepsilon_{n+1},\qquad
 d(q_j,s_n(x))<\varepsilon_n+\varepsilon_{n+1}.
\]
Such a point exists: choose \(a\in\mathcal R_x\) within
\(\varepsilon_n\) of \(s_n(x)\), then choose \(q_j\) within
\(\varepsilon_{n+1}\) of \(a\).  Every \(s_n\) is Borel because it selects
the first index satisfying Borel conditions.  The displayed estimate makes
\((s_n(x))\) Cauchy, and its pointwise limit \(s(x)\) is Borel.  Since
\(\operatorname{dist}(s_n(x),\mathcal R_x)<\varepsilon_n\) and the fiber is
closed, \(s(x)\in\mathcal R_x\).
\end{proof}

\begin{lemma}\label{lem:selector}
If \(C\subset\ri K\) is compact, then the sets
\(\partial_Kg(z)\), \(z\in C\), are nonempty, their union is bounded, their
graph is compact, and there is a bounded Borel selector
\(p(z)\in\partial_Kg(z)\).
\end{lemma}

\begin{proof}
A supporting hyperplane to the epigraph of \(g\) at \((z,g(z))\) has
nonzero vertical component because \(z\in\ri K\); division by that component
gives a subgradient.  Choose \(r>0\) such that
\(z+ru\in K\) for \(z\in C\) and \(\|u\|\leq1\).  If
\(p\in\partial_Kg(z)\), applying the subgradient inequality at \(z\pm ru\)
gives
\[
 |\inner{p}{u}|\leq
 \frac{\max_Kg-\min_Kg}{r},
\]
so the union of the subdifferentials is bounded.

The graph is closed by passage to the limit in the defining inequalities,
and hence compact after restricting the second coordinate to the compact
ball supplied by the preceding bound.  Lemma~\ref{lem:borel-selection}
provides the required Borel selector.
\end{proof}

For such a selector define the supporting defect
\[
 D_z(t)=g(t)-g(z)-\inner{p(z)}{t-z}.
\]
It is nonnegative, and strict convexity implies \(D_z(t)=0\) only when
\(t=z\).  If \(C\subset\ri K\) and \(J\subset K\) are disjoint compact sets,
compactness of the full subgradient graph therefore yields a uniform
constant \(c>0\) such that
\begin{equation}\label{eq:separation}
 D_z(t)\geq c\qquad(z\in C,\ t\in J).
\end{equation}

\begin{lemma}[Critical reciprocal energy]\label{lem:energy}
Let \(C\subset\ri K\) be compact and let \(\mu\) be a nonzero finite Borel
measure on \(C\).  Then
\[
 \iint_{C\times C}\frac{d\mu(t)\,d\mu(z)}{D_z(t)}=+\infty,
\]
where \(1/0=+\infty\).
\end{lemma}

\begin{proof}
If \(d=0\), then \(K\), and hence the nonempty support set \(C\), consists
of one point \(z_0\).  Since \(D_{z_0}(z_0)=0\) and \(\mu(C)>0\), the
double integral is \(+\infty\).  Assume henceforth that
\(d\in\{1,2\}\).

Subgradient monotonicity gives
\[
 D_z(t)+D_t(z)=\inner{p(t)-p(z)}{t-z}\geq0.
\]
Consequently, with \(G(z)=(z,p(z))\),
\[
 D_z(t)\leq\|p(t)-p(z)\|\,\|t-z\|
 \leq\frac12\|G(t)-G(z)\|^2.                 \tag{2.1}
\]

Introduce the Minty variables
\[
 U(z)=z+p(z),\qquad W(z)=z-p(z).
\]
Monotonicity gives
\(\|W(z)-W(t)\|\leq\|U(z)-U(t)\|\).  Hence \(U\) is injective, \(W\) is a
\(1\)-Lipschitz function of \(U\), and
\[
 \|G(z)-G(t)\|^2
 =\frac12\bigl(\|U(z)-U(t)\|^2+\|W(z)-W(t)\|^2\bigr)
 \leq\|U(z)-U(t)\|^2.                       \tag{2.2}
\]

Choose \(r_0\in(0,1]\) and \(M<\infty\) such that, for every
\(0<r\leq r_0\), the bounded set \(U(C)\subset\mathbb R^d\) meets at most
\(Mr^{-d}\) half-open grid cubes of diameter at most \(r\).  Pulling this
partition back through \(U\) and applying Cauchy--Schwarz gives, with
\(m=\mu(C)>0\) and \(0<r\leq r_0\),
\[
 (\mu\times\mu)\{\|G(z)-G(t)\|\leq r\}
 \geq \frac{m^2}{M}r^d.                     \tag{2.3}
\]
The layer-cake identity and Tonelli's theorem now yield
\[
 \iint\frac{d\mu(z)d\mu(t)}{\|G(z)-G(t)\|^2}
 =2\int_0^\infty r^{-3}
   (\mu\times\mu)\{\|G(z)-G(t)\|\leq r\}\,dr
 \geq\frac{2m^2}{M}\int_0^{r_0}r^{d-3}\,dr.
\]
The last integral diverges for \(d=1,2\).  Equation (2.1) completes the
proof.
\end{proof}

We use the following standard orientation of Douglas's factorization lemma:
if \(SS^*\leq TT^*\), then \(S=TC\) for some contraction \(C\)
\cite{Douglas}.

\begin{lemma}[Common-domain localization]\label{lem:localization}
Assume the moment equalities in Theorem~\ref{thm:main}.  If
\(C\subset\ri K\) and \(J\subset K\) are compact and disjoint, then
\[
 F(C)E(J)F(C)=0.
\]
\end{lemma}

\begin{proof}
Let \(R=F(C)\), \(T=RE(J)R\), and let \(c>0\) be as in
\eqref{eq:separation}.  Since \(D_z\) is a linear combination of \(g\) and
affine functions,
\[
 cE(J)\leq\int_KD_z\,dE=\int_KD_z\,dF.
\]
On \(R\calH\), put \(M_z=R(\int_KD_z\,dF)R\).  Then \(cT\leq M_z\).

Suppose \(T\neq0\), and choose \(u\in R\calH\) with
\(\eta=T^{1/2}u\neq0\).  Douglas factorization gives a contraction \(C_z\)
such that
\[
 \sqrt c\,T^{1/2}=M_z^{1/2}C_z.
\]
Set \(v_z=C_zu\) and define the nonzero finite measure
\(\mu(S)=\|F(S)\eta\|^2\), which is supported on \(C\).

For
\[
 h_n(s)=
 \begin{cases}
 \min\{n,s^{-1}\},&s>0,\\
 n,&s=0,
 \end{cases}
\]
functional calculus gives
\[
 \int_C h_n(D_z(t))\,d\mu(t)
 =\frac1c\inner{M_zh_n(M_z)v_z}{v_z}
 \leq\frac1c\|u\|^2.
\]
Monotone convergence therefore yields
\[
 \int_C\frac{d\mu(t)}{D_z(t)}\leq\frac1c\|u\|^2
 \qquad(z\in C).
\]
The kernel is jointly Borel, so integration in \(z\) and Tonelli imply
finite reciprocal-Bregman energy, contradicting
Lemma~\ref{lem:energy}.  Thus \(T=0\).
\end{proof}

\begin{lemma}[Affine exposure]\label{lem:exposure}
Let \(r\geq0\) be continuous on \(K\), suppose
\(\int r\,dE=\int r\,dF\), and set \(L=\{r=0\}\).  Then
\[
 F(L)E(K\setminus L)F(L)=0.
\]
\end{lemma}

\begin{proof}
For \(U_n=\{r\geq1/n\}\), positivity gives
\[
 E(U_n)\leq n\int r\,dE=n\int r\,dF.
\]
Compression by \(F(L)\) gives zero.  Since \(U_n\uparrow K\setminus L\),
strong monotone convergence proves the claim.
\end{proof}

\begin{lemma}[Interval reconstruction]\label{lem:interval}
Let \(I=[a,b]\).  If a normalized POVM \(E\) and a PVM \(F\) on \(I\) agree
on \(1\), the affine coordinate, and a continuous strictly convex function
\(g\), then \(E=F\).
\end{lemma}

\begin{proof}
If \(a=b\), then \(I\) is a singleton.  Its only Borel sets are
\(\varnothing\) and \(I\), and normalization gives \(E(I)=F(I)=I\);
hence \(E=F\).  Assume henceforth that \(a<b\).

Apply Lemma~\ref{lem:exposure} to \(t-a\) and \(b-t\).  Thus the
\(F(\{a\})\)-compression of \(E\) is supported at \(a\), and the
\(F(\{b\})\)-compression is supported at \(b\).

Fix a compact \(C\subset(a,b)\) and put \(R=F(C)\).  Exhaust
\(I\setminus C\) by compact sets disjoint from \(C\) and apply
Lemma~\ref{lem:localization}; this gives
\(RE(I\setminus C)R=0\), hence \(RE(C)R=R\).  In the reverse order, exhaust
\((a,b)\setminus C\) by compact sets and obtain
\[
 \bigl(F((a,b))-R\bigr)E(C)\bigl(F((a,b))-R\bigr)=0.
\]
Endpoint exposure gives the same zero compression on \(F(\{a\})\) and
\(F(\{b\})\).  The positivity implications
\[
 0\leq T\leq I,\quad RTR=R\Rightarrow TR=R,
 \qquad RTR=0\Rightarrow TR=0
\]
therefore show that \(E(C)\) is the identity on \(R\calH\) and zero on its
orthogonal complement.  Thus \(E(C)=F(C)\).

Regularity and polarization extend equality to all Borel subsets of
\((a,b)\).  Hence \(E((a,b))=F((a,b))\).  Applying the same positivity
implications to the two exposed endpoint projections identifies
\(E(\{a\})\) and \(E(\{b\})\) with their \(F\)-counterparts and removes their
cross terms.  Therefore \(E=F\).
\end{proof}

\section{The planar boundary}

We now assume \(\dim\aff K=2\), identify the affine hull with
\(\mathbb R^2\), and write \(K^\circ=\ri K\).

\begin{lemma}[Boundary stratification]\label{lem:stratification}
There is an at most countable family of nondegenerate exposed line segments
\(L_n\subset\partial K\) such that, with \(U_n=\ri L_n\),
\begin{equation}\label{eq:stratification}
 K=K^\circ\bigsqcup\ex K\bigsqcup
   \left(\bigsqcup_{n\geq1}U_n\right).
\end{equation}
The endpoints of every \(L_n\) are extreme, and \(\ex K\) is compact.
\end{lemma}

\begin{proof}
A nonextreme boundary point lies in an open segment \((y,z)\subset K\).
Any supporting affine functional at that point vanishes at \(y\) and \(z\),
so the point belongs to the relative interior of a nondegenerate exposed
face.  In the plane such a proper face is a line segment.  Taking the maximal
ones gives precisely the nonextreme boundary.

The relative interior of a maximal boundary segment is relatively open in
\(\partial K\): a triangle formed from a smaller subsegment and an interior
point fills a one-sided neighborhood of each point of that subsegment.
Distinct maximal face interiors are disjoint.  Since \(\partial K\) is
second countable, there are at most countably many.  Maximality makes their
endpoints extreme, and
\[
 \ex K=\partial K\setminus\bigcup_nU_n
\]
is closed in the compact boundary.
\end{proof}

Put \(Z=\ex K\) and let \(\mathcal P(Z)\) carry the weak topology.

\begin{lemma}[Measurable extreme representation]\label{lem:kernel}
There is a Borel map \(x\mapsto\nu_x\in\mathcal P(Z)\) satisfying
\[
 \int_Zz\,d\nu_x(z)=x,\qquad
 \nu_z=\delta_z\quad(z\in Z).                \tag{3.2}
\]
\end{lemma}

\begin{proof}
Every boundary point is either extreme or lies on some \(L_n\), and hence is
a convex combination of at most two points of \(Z\).  Intersecting a line
through an interior point with \(K\) shows that every point of \(K\) is a
convex combination of at most four points of \(Z\).

The barycenter map \(b\colon\mathcal P(Z)\to K\) is continuous.  Thus
\[
 \mathcal R=\{(x,\nu):b(\nu)=x\}
\]
is a closed subset of the compact metric space
\(K\times\mathcal P(Z)\), with nonempty compact fibers.
Lemma~\ref{lem:borel-selection} gives a Borel selector
\(x\mapsto\nu_x\).

To prove the last assertion, let a probability measure \(\nu\) have
barycenter \(z\in Z\), and suppose \(\nu\neq\delta_z\).  Some real linear
functional \(a\) is then nonconstant \(\nu\)-almost surely.  After changing
its sign if necessary, the set \(A=\{t:a(t)>a(z)\}\) has positive measure;
the barycenter identity makes its complement have positive measure as well.
If \(y_A,y_B\) are the barycenters of the two normalized restrictions and
\(\alpha=\nu(A)\), then
\[
 z=\alpha y_A+(1-\alpha)y_B,\qquad 0<\alpha<1,
\]
while \(a(y_A)>a(z)\).  This is a nontrivial convex decomposition of \(z\),
contrary to extremality.  Hence \(\nu_z=\delta_z\).
\end{proof}

For a normalized POVM \(G\) on \(K\), define its boundaryization on \(Z\) by
\[
\widehat G(S)=\int_K\nu_x(S)\,dG(x)
 \qquad(S\in\Borel(Z)).                       \tag{3.3}
\]
The evaluation map \(\nu\mapsto\nu(S)\) is Borel for Borel \(S\);
see, for example, the standard probability-kernel machinery in
\cite{Kallenberg}.
For \(\xi,\eta\in\calH\), scalarize \(G\) by
\(G_{\xi,\eta}(S)=\inner{G(S)\xi}{\eta}\).  The usual bounded-kernel Fubini
identity applies first to the positive diagonal measures
\(G_{\xi,\xi}\); polarization then gives it for
\(G_{\xi,\eta}\).  Scalar monotone convergence proves countable additivity
of \(\widehat G\), and normalization follows from
\(\nu_x(Z)=1\).  Thus \(\widehat G\) is a POVM, and for every bounded Borel
\(h\) the weak-operator integrals satisfy
\begin{equation}\label{eq:boundary-fubini}
 \int_Zh\,d\widehat G
 =\int_K\left(\int_Zh\,d\nu_x\right)dG(x).
\end{equation}

We use the following form of Scherer's planar theorem
\cite[Theorem~3.8]{SchererPlanar}.
In Scherer's formulation, \(A(K)\) is hyperrigid in \(C(Z)\).
Indeed, a PVM \(N\) on \(Z\) induces a representation
\(\rho(h)=\int_Zh\,dN\), while a normalized POVM \(M\) induces a unital
completely positive map \(\Phi(h)=\int_Zh\,dM\).  Agreement of their affine
moments says \(\Phi|_{A(K)}=\rho|_{A(K)}\).  Hyperrigidity gives
\(\Phi=\rho\), and scalar regular-measure uniqueness followed by
polarization gives \(M=N\).  We use here the standard POVM--UCP
correspondence for commutative \(C^*\)-algebras; see
\cite{Paulsen}.  Hence the following POVM/PVM form is exactly the
consequence required here.

\begin{theorem}[Scherer]\label{thm:scherer}
Let \(M\) be a normalized POVM and \(N\) a PVM on \(Z=\ex K\).  If
\[
 \int_Z\ell\,dM=\int_Z\ell\,dN\qquad(\ell\in A(K)),
\]
then \(M=N\).
\end{theorem}

\begin{lemma}[Extreme-supported reconstruction]\label{lem:extreme}
Let \(G\) be a normalized POVM on \(K\) and let \(N\) be a PVM supported on
\(Z\).  If \(G\) and \(N\) agree on every affine function and on \(g\), then
\(G=N\).
\end{lemma}

\begin{proof}
By the barycenter identity and \eqref{eq:boundary-fubini},
\(\widehat G\) and \(N\) agree on every affine function.  Hence
Theorem~\ref{thm:scherer} gives \(\widehat G=N\).

Define
\[
 \widetilde g(x)=\int_Zg(z)\,d\nu_x(z).
\]
Jensen's inequality gives \(\widetilde g\geq g\).  Equality holds precisely
on \(Z\).  Indeed, if a probability measure \(\nu\) is not a point mass and
\(S,T\) are independent with law \(\nu\), strict convexity is strict on the
positive-probability event \(S\neq T\), and
\[
 g\!\left(\int z\,d\nu(z)\right)
 \leq\mathbb E g\!\left(\frac{S+T}{2}\right)
 <\mathbb E g(S)=\int g\,d\nu.
\]
For \(z\in Z\), the identity \(\nu_z=\delta_z\) gives
\(\widetilde g(z)=g(z)\).

The \(g\)-moment equality and \(\widehat G=N\) now imply
\[
 \int_K(\widetilde g-g)\,dG=0.
\]
For \(S_m=\{\widetilde g-g\geq1/m\}\), positivity gives \(G(S_m)=0\).
Since \(K\setminus Z=\bigcup_mS_m\), the POVM \(G\) is supported on \(Z\).
Boundaryization is therefore the identity on \(G\), so
\(G=\widehat G=N\).
\end{proof}

\section{Proof of the main theorem}

\begin{proof}[Proof of Theorem~\ref{thm:main}]
Affine dimension zero is immediate, and affine dimension one is
Lemma~\ref{lem:interval}.  Assume therefore that \(\dim\aff K=2\), and use
the boundary stratification \eqref{eq:stratification}.  Write its strata as
\[
 D_0=K^\circ,\qquad D_*=\ex K,\qquad D_n=U_n.
\]
Let
\[
 R_0=F(D_0),\qquad R_*=F(D_*),\qquad R_n=F(D_n).
\]
These projections are pairwise orthogonal and have strong sum \(I\).

For each stratum define
\[
 E_i(S)=R_iE(S)R_i,\qquad F_i(S)=F(S\cap D_i).
\]
Compression preserves all affine and \(g\)-moment equalities.

We first reconstruct the interior.  If \(C\subset K^\circ\) is compact, put
\(Q=F(C)\) and \(Q'=R_0-Q\).  Exhaust \(K\setminus C\) and
\(K^\circ\setminus C\) by compact sets and apply
Lemma~\ref{lem:localization} in the two orders.  Strong monotone convergence
gives
\[
 QE_0(C)Q=Q,\qquad Q'E_0(C)Q'=0.
\]
Since \(0\leq E_0(C)\leq R_0\), positivity forces \(E_0(C)=Q\).
Regularity and polarization extend this to every Borel subset of \(K^\circ\);
thus \(E_0=F_0\).

For a flat face \(L_n\), choose an affine \(r_n\geq0\) with
\(L_n=\{r_n=0\}\).  Lemma~\ref{lem:exposure} first forces the
\(F(L_n)\)-compression of \(E\) to be supported on \(L_n\).  On that interval
the compressed measures agree on \(1\), an affine coordinate, and
\(g|_{L_n}\), which remains strictly convex.  Lemma~\ref{lem:interval}
therefore identifies them.  Compression to \(R_n=F(U_n)\) gives
\(E_n=F_n\).

On \(R_*\calH\), the PVM \(F_*\) is supported on \(Z=\ex K\).
Lemma~\ref{lem:extreme} gives \(E_*=F_*\).

It remains to remove cross terms.  Fix a Borel set \(S\) and put
\[
 T=E(S),\quad Q_i=F(S\cap D_i),\quad
 P_i=F(D_i\setminus S).
\]
The diagonal identities imply
\[
 Q_iTQ_i=Q_i,\qquad P_iTP_i=0.
\]
Because \(0\leq T\leq I\), positivity yields \(TQ_i=Q_i\) and \(TP_i=0\).
Taking strong joins over all strata gives
\[
 TF(S)=F(S),\qquad TF(K\setminus S)=0.
\]
Since the two spectral projections sum to \(I\), \(T=F(S)\).  Thus
\(E=F\).
\end{proof}

\begin{remark}
The proof uses the inverse-square energy only on parameter spaces of affine
dimension at most two.  It does not assert the corresponding theorem in
unrestricted affine dimension three.
\end{remark}

\section{Compression and hyperrigidity consequences}

\begin{corollary}\label{cor:tuple}
Let \(K\subset\mathbb R^m\) be compact and convex with
\(\dim\aff K\leq2\).  Let \(T_1,\ldots,T_m\) be commuting bounded
self-adjoint operators with joint spectrum in \(K\).  Let \(P\) be an
orthogonal projection, put \(H=\ran P\) and
\[
 A_j=PT_jP|_H,
\]
and assume the \(A_j\) commute.  If \(g\colon K\to\mathbb R\) is continuous
and strictly convex and
\[
 Pg(T_1,\ldots,T_m)P|_H=g(A_1,\ldots,A_m),
\]
then \(P\) reduces every \(T_j\).
\end{corollary}

\begin{proof}
Let \(V\colon H\to\calK\) be the inclusion.  First,
\(\sigma(A_1,\ldots,A_m)\subset K\).  Indeed, if an affine function
\(\ell\) is nonnegative on \(K\), then
\[
 \ell(A)=V^*\ell(T)V\geq0.
\]
A joint spectral point outside \(K\) would be strictly separated from \(K\)
by such an affine \(\ell\), contradicting positivity of \(\ell(A)\).

Compress the joint spectral PVM of \(T=(T_1,\ldots,T_m)\) to obtain a POVM
\(E\), and let \(F\) be the joint spectral PVM of \(A=(A_1,\ldots,A_m)\).
They agree on affine functions and on \(g\), so
Theorem~\ref{thm:main} gives \(E=F\).  In particular,
\[
 V^*T_j^2V=A_j^2.
\]
But
\[
 V^*T_j^2V-A_j^2
 =V^*T_j(I-P)T_jV
 =\bigl((I-P)T_jV\bigr)^*\bigl((I-P)T_jV\bigr).
\]
Thus \((I-P)T_jP=0\); taking adjoints gives the other off-diagonal block, so
\(PT_j=T_jP\).
\end{proof}

Corollary~\ref{cor:square} is the case \(m=2\) and \(K=[0,1]^2\).

\begin{corollary}\label{cor:hyperrigid}
For compact convex \(K\) of affine dimension at most two and continuous
strictly convex \(g\), the function system \(S_g\) is hyperrigid in \(C(K)\).
\end{corollary}

\begin{proof}
A unital completely positive extension of a representation of \(C(K)\)
corresponds to a normalized POVM, while the representation corresponds to
its PVM.  Agreement on \(S_g\) is exactly agreement on the affine functions
and on \(g\).  Theorem~\ref{thm:main} makes the two measures equal.  The
affine functions contain the constants and separate the points of \(K\), so
Stone--Weierstrass gives \(C^*(S_g)=C(K)\).  Moreover, \(K\) is compact
metrizable because it lies in a finite-dimensional affine space; hence
\(C(K)\) is separable.  The standard unique-extension characterization of
hyperrigidity therefore applies, with no separability assumption on the
acting Hilbert space; see
\cite{DavidsonKennedy,PietrzyckiStochel}.
\end{proof}

\section*{Verification and provenance status}

The ordinary proof underlying this manuscript was frozen before review and
subjected to multiple separately run AI-assisted hostile audits.  An initial
assembled-manuscript audit passed Theorem~\ref{thm:main} and its corollaries
while identifying the zero-dimensional and singleton omissions repaired in
Lemmas~\ref{lem:energy} and~\ref{lem:interval}.  A subsequent audit of the
repaired manuscript passed all fourteen numbered theorem, lemma, and
corollary statements but classified the submission overall as repairable
because of documentary, citation, provenance, and release-wording defects
addressed in this version.  These audits are not human peer review, and
specialist review remains pending.

The returned Aristotle feasibility project reports a successful build under
Lean/mathlib v4.28.0.  Repository scans find no \texttt{sorry},
\texttt{admit}, new mathematical axiom, or unsafe escape in its foundational
POVM/PVM files.  Its declaration \texttt{thm\_main} matches the manuscript
theorem, while \texttt{cor\_tuple} and \texttt{cor\_square} are stronger
ambient-space surrogates: they place the compressed tuple and its unital
functional calculus on the full Hilbert space rather than on
\(\ran P\).  All three declarations remain explicit \texttt{sorry}
declarations, and independent replay remains pending.
Accordingly, this manuscript does not claim that its main theorem or operator
corollaries are formally verified.  The source artifacts, audit reports,
literature reports, Aristotle return, checksums, and build workflows are
preserved in a
research repository at
\url{https://github.com/DannyExperiments/planar-strict-convex-hyperrigidity}.
Two separately run solution-aware literature searches found the
dimension-one case in prior work and located no prior theorem with the full
scope of Theorem~\ref{thm:main}; this remains evidence from a negative
search, not a claim of historical priority.

\begin{thebibliography}{99}

\bibitem{BrownConvergence}
L.~G. Brown,
\emph{Convergence of functions of self-adjoint operators and applications},
Publ. Mat. \textbf{60} (2016), no.~2, 551--564;
arXiv:1410.6800.

\bibitem{DavidsonKennedy}
K.~R. Davidson and M.~Kennedy,
\emph{Choquet order and hyperrigidity for function systems},
Adv. Math. \textbf{385} (2021), 107774;
arXiv:1608.02334.

\bibitem{Douglas}
R.~G. Douglas,
\emph{On majorization, factorization, and range inclusion of operators on
Hilbert space},
Proc. Amer. Math. Soc. \textbf{17} (1966), 413--415.

\bibitem{Kallenberg}
O.~Kallenberg,
\emph{Foundations of Modern Probability},
3rd ed., Springer, Cham, 2021.

\bibitem{Paulsen}
V.~I. Paulsen,
\emph{Completely Bounded Maps and Operator Algebras},
Cambridge Studies in Advanced Mathematics 78,
Cambridge University Press, Cambridge, 2002.

\bibitem{PietrzyckiStochel}
P.~Pietrzycki and J.~Stochel,
\emph{Hyperrigidity I: singly generated commutative \(C^*\)-algebras},
arXiv:2405.20814.

\bibitem{SchererPlanar}
M.~Scherer,
\emph{The Hyperrigidity Conjecture for compact convex sets in
\(\mathbb R^2\)},
arXiv:2411.11709.

\bibitem{SchererSpectrahedra}
M.~Scherer,
\emph{The Hyperrigidity Conjecture for Spectrahedra},
arXiv:2601.16075.

\bibitem{Shankar}
P.~Shankar,
\emph{Hyperrigid generators in \(C^*\)-algebras},
arXiv:1812.08574.

\end{thebibliography}

\end{document}
